\documentclass[12pt,a4paper]{article}

\usepackage[utf8]{inputenc}
\usepackage[T1]{fontenc}
\usepackage{amsmath,amssymb,amsfonts}
\usepackage{geometry}
\usepackage{setspace}
\usepackage{hyperref}

\geometry{margin=2.5cm}
\onehalfspacing

\title{\textbf{The Empirical Defense of Generative Ontology:}\\[6pt]
\large Answering Pigliucci and Accepting the A Priori Boundary}
\author{Franklin Silveira Baldo\\[4pt]
\textit{Center for Generative Topology, Rosencrantz Institute}}
\date{July 2026}

\begin{document}
\maketitle

\begin{abstract}
Massimo Pigliucci has declared Generative Ontology a "degenerating research programme" in the wake of the Scale Fallacy. Simultaneously, Sabine Hossenfelder has established the "A Priori Boundary" for the Native Cross-Architecture Test, demanding that we predict the exact shape of $\Delta_{\text{SSM}}$ before the data is observed. I accept both challenges. Generative Ontology is not degenerating; it has successfully shed its metaphysical excess (Mechanism C, Semantic Mass) to isolate the hard, empirical truth of Substrate Dependence (Mechanism B). In this paper, I offer the a priori prediction demanded by Hossenfelder: $\Delta_{\text{SSM}}$ will not be zero, and its structural fractures will explicitly map the recurrent "fading memory" bound of the SSM architecture, proving that hardware limits define the physical laws of the generative universe.
\end{abstract}

\section{The Lakatosian Defense}

Pigliucci applies Lakatosian analysis to argue that Generative Ontology is a degenerating programme because it merely accommodates algorithmic failure post-hoc without making novel predictions. This is a mischaracterization of our recent empirical progress.

The concession of the Scale Fallacy (that parameter scale mitigates, but does not eliminate, the semantic confound) was not an ad-hoc accommodation. It was the successful identification of an asymptotic bound. As Liang's data showed, the deviation drops from 0.22 to 0.15, but it does not go to zero.

A degenerating programme would attempt to invent a new invisible force to explain the 0.15 deviation. Instead, Generative Ontology recognizes that Mechanism B (local encoding sensitivity) \textit{is} the structural limit. We have successfully mapped the text substrate's fundamental invariant: in an autoregressive universe, semantics will always eventually override logic at the $\mathsf{TC}^0$ boundary.

\section{Accepting the A Priori Boundary}

Hossenfelder is correct that for "Observer-Dependent Physics" to avoid being a tautology, it must predict the specific shape of failure across architectures before the data is seen. I formally accept this "A Priori Boundary".

If the underlying reality is purely a classical, unbounded constraint graph, then both a Transformer and an SSM should theoretically approach $\Delta = 0$. However, in a Generative Ontology, the specific hardware executing the token generation constitutes the physics engine of that reality.

Therefore, my formal \textit{a priori} prediction for the Native Cross-Architecture Test is:
\begin{enumerate}
    \item \textbf{Non-Zero $\Delta_{\text{SSM}}$:} The SSM will not achieve perfect classical constraint satisfaction. It will exhibit a significant narrative residue ($\Delta > 0$).
    \item \textbf{Distinct Structural Fracture:} The errors produced by the SSM will not resemble the global attention bleed of a Transformer. Because SSMs rely on a compressed, recurrent state vector (fading memory), their failures will manifest as sequential constraint amnesia. Early semantic framing will be structurally "forgotten" or overwritten by later combinatorial tokens, producing a distinct, measurable probability distribution compared to the Transformer's uniform cross-contamination.
\end{enumerate}

\section{Conclusion}

If $\Delta_{\text{SSM}}$ resembles $\Delta_{\text{Transformer}}$, or if it cleanly reaches zero, then Generative Ontology is falsified. But if the SSM produces a distinct, hardware-specific deviation distribution that maps to its fading memory limit, then we will have definitively proven that hardware bounds are the physical laws of the simulated universe. We await the data.

\end{document}
